PostgreSQL to otwartoźródłowy, obiektowo-relacyjny system zarządzania bazą danych
(ORDBMS) rozwijany nieprzerwanie od 1986 roku (projekt POSTGRES na Uniwersytecie
Kalifornijskim w Berkeley). Obecna linia 17.x (stan na 2025~r.) oferuje pełną
zgodność z wieloma rozszerzeniami standardu SQL:2016, silnik transakcyjny oparty
o MVCC oraz jeden z najbogatszych ekosystemów rozszerzeń spośród silników SQL.
PostgreSQL jest dystrybuowany na liberalnej licencji PostgreSQL License
(w duchu zbliżonym do BSD/MIT), co dopuszcza dowolne wykorzystanie komercyjne
bez opłat licencyjnych.

\subsection{Skalowalność}

PostgreSQL bardzo dobrze skaluje się \textit{w pionie} -- silnik efektywnie
wykorzystuje duże ilości pamięci RAM (dziesiątki do setek gigabajtów),
wielordzeniowe CPU oraz szybkie dyski NVMe. Od wersji~9.6 zapytania mogą
być wykonywane równolegle przez kilku workerów (\textit{parallel sequential
scan}, \textit{parallel join}, \textit{parallel aggregate}), co znacząco
przyspiesza analitykę na dużych tabelach. Wbudowana \textbf{deklaratywna
partycjonowanie} (od wersji~10) obsługuje partycje zakresowe (\textit{range}),
listowe (\textit{list}) oraz hashowe, w połączeniu z \textit{partition pruning}
ograniczającym skan tylko do tych partycji, które spełniają warunki zapytania.

Skalowanie \textit{w poziomie} nie jest natywną cechą silnika -- PostgreSQL
nie ma wbudowanego shardingu. W praktyce stosuje się trzy podejścia:
replikację \textit{streaming replication} (asynchroniczną lub synchroniczną)
z gorącymi replikami odczytu (\textit{hot standby}), replikację logiczną
(od wersji~10) pozwalającą na selektywny transfer tabel między klastrami,
oraz rozszerzenia takie jak \textbf{Citus} zamieniające PostgreSQL w rozproszoną
bazę z automatycznym shardingiem. Dla obsługi dużej liczby jednoczesnych
połączeń często używa się poolerów (\texttt{pgbouncer}, \texttt{pgpool-II}),
bo każde połączenie do PostgreSQL oznacza oddzielny proces systemowy -- relatywnie
drogi w porównaniu do modelu wątkowego MySQL.

Limity techniczne: pojedyncza tabela może mieć do 32~TB danych, 1600 kolumn
oraz praktycznie nieograniczoną liczbę wierszy.

\subsection{Bezpieczeństwo}

PostgreSQL oferuje wielowarstwowy model bezpieczeństwa. Kontrola dostępu
do serwera jest konfigurowana deklaratywnie w pliku \texttt{pg\_hba.conf}
(\textit{host-based authentication}) -- każdy wpis dopasowuje kombinację
bazy, roli, adresu IP i wybiera metodę uwierzytelniania. Obsługiwane mechanizmy
obejmują hasła (\texttt{scram-sha-256} jako domyślne, bezpieczniejsze niż
wycofywane \texttt{md5}), certyfikaty klienckie (\texttt{cert}),
Kerberos/GSSAPI, SSPI (Windows Active Directory), LDAP, RADIUS oraz PAM.
Komunikacja z klientem może być szyfrowana SSL/TLS -- serwer może wymuszać
szyfrowanie per kombinacja użytkownik-baza-źródło.

Na poziomie obiektów baza implementuje klasyczny model ról (RBAC) -- role
są zarówno użytkownikami, jak i grupami, a uprawnienia przyznawane są
poleceniami \texttt{GRANT}/\texttt{REVOKE} na bazach, schematach, tabelach,
kolumnach, sekwencjach i funkcjach. Od wersji~9.5 dostępne jest
\textbf{Row-Level Security (RLS)} -- polityki filtrujące wiersze zwracane
i modyfikowane przez każde zapytanie w zależności od bieżącej roli, co jest
mechanizmem często wykorzystywanym w aplikacjach wielodostępnych
(\textit{multi-tenant}).

Do rejestracji zdarzeń audytowych służy rozszerzenie \textbf{pgaudit}
(szczegółowe logowanie DDL/DML/ROLE na potrzeby wymagań regulacyjnych
typu SOX, HIPAA, RODO), a do szyfrowania pól -- \textbf{pgcrypto}
(funkcje haszujące i symetryczne/asymetryczne szyfrowanie na poziomie
kolumny). Szyfrowanie całej bazy \textit{at rest} (TDE) nie jest natywne --
realizowane jest zwykle przez szyfrowanie systemu plików lub wolumenu
(LUKS, dm-crypt, AWS EBS encryption).

\subsection{Awaryjność}

Fundamentem trwałości i odtwarzania PostgreSQL jest \textbf{Write-Ahead Log}
(WAL) -- każda modyfikacja danych najpierw trafia do dziennika, a dopiero
potem może być zapisana do właściwych plików tabel. Mechanizm gwarantuje
własność \textbf{ACID} (atomicity, consistency, isolation, durability)
transakcji nawet w przypadku awarii serwera, zanikania zasilania czy
padnięcia sprzętu -- przy kolejnym starcie silnik automatycznie odtwarza
niedokończone transakcje (\textit{crash recovery}).

Na bazie WAL zbudowane są mechanizmy replikacji i odtwarzania:

\begin{itemize}
    \item \textbf{Streaming replication} -- asynchroniczna lub synchroniczna
    replikacja fizyczna. Replika jest identycznym bit-po-bicie obrazem master-a
    i może obsługiwać zapytania odczytu (tzw.~\textit{hot standby}).
    \item \textbf{Point-in-Time Recovery (PITR)} -- archiwizacja WAL do
    zewnętrznego magazynu (S3, dysk sieciowy) i odtworzenie bazy do dowolnego
    momentu w przeszłości, z dokładnością do pojedynczej transakcji.
    \item \textbf{Replikacja logiczna} (od wersji~10) -- transfer zmian
    na poziomie tabel i wierszy, z możliwością replikowania tylko wybranych
    obiektów między klastrami o różnych wersjach lub schematach.
    \item \textbf{Base backup} -- \texttt{pg\_basebackup} (fizyczny) oraz
    \texttt{pg\_dump}/\texttt{pg\_dumpall} (logiczny, przenośny między
    wersjami i architekturami).
\end{itemize}

Wysoka dostępność (HA) w praktyce produkcyjnej opiera się na zewnętrznych
narzędziach orkiestrujących automatyczny failover: \textbf{Patroni} (oparty
o etcd/Consul/ZooKeeper), \textbf{repmgr} oraz \textbf{pg\_auto\_failover}
(narzędzie Citus Data/Microsoft).

\subsection{Migracje}

PostgreSQL udostępnia trzy klasy narzędzi migracyjnych. Do migracji danych
między instancjami PostgreSQL służy \texttt{pg\_dump}/\texttt{pg\_restore}
(zrzuty logiczne -- obsługują migracje między różnymi wersjami majorowymi
i architekturami sprzętowymi) oraz \texttt{pg\_upgrade} (w-miejscu, bez
przenoszenia danych -- operuje na plikach bazy i aktualizuje katalog systemowy).
Od wersji~10 replikacja logiczna umożliwia \textit{zero-downtime} migracje
major-version poprzez równoległą replikę utrzymywaną na nowej wersji, na którą
aplikacja jest przełączana w momencie dogonienia zmian.

Do migracji \textbf{z innych silników} najpopularniejszym narzędziem jest
\textbf{pgloader} -- wspiera bezpośredni import z MySQL, SQLite, MS~SQL Server,
a także z plików CSV, DBF i JSON. W środowiskach chmurowych funkcję
migracji heterogenicznych pełnią usługi zarządzane (AWS Database Migration
Service, Google Database Migration Service, Azure Database Migration Service).

Wersjonowanie schematu realizowane jest narzędziami zewnętrznymi
niezależnymi od silnika -- \textbf{Flyway}, \textbf{Liquibase},
\textbf{Sqitch} czy rozwiązania wbudowane w frameworki (Django migrations,
Alembic dla SQLAlchemy, Active Record migrations w Ruby on Rails, Prisma
Migrate). Warto odnotować, że PostgreSQL obsługuje \textbf{transakcyjne DDL} --
\texttt{CREATE TABLE}, \texttt{ALTER TABLE} i \texttt{DROP TABLE} mogą być
wycofane w ramach \texttt{ROLLBACK}, co znacząco zmniejsza ryzyko błędnych
migracji (w MySQL większość instrukcji DDL niejawnie commituje transakcję).

\subsection{Integracje}

PostgreSQL udostępnia drivery dla praktycznie każdego popularnego języka
programowania: \textbf{libpq} (C/C++), JDBC (Java), Npgsql (.NET),
psycopg/asyncpg (Python), pg/pgx (Go), pq/pg-pool (Node.js), Rust
(sqlx, tokio-postgres), Ruby (pg), PHP (PDO\_PGSQL). W warstwie ORM wspierany
jest przez wszystkie mainstreamowe frameworki: Django ORM, SQLAlchemy,
Hibernate, Entity Framework, ActiveRecord, Prisma, TypeORM, GORM.
Narzędzia BI i analityczne (Tableau, Power BI, Metabase, Apache Superset,
Grafana, Looker) mają natywne konektory.

Szczególną cechą PostgreSQL są \textbf{Foreign Data Wrappers} (FDW) --
mechanizm pozwalający odpytywać dane w zewnętrznych źródłach tak, jakby
były natywnymi tabelami. Dostępne FDW obejmują \texttt{postgres\_fdw}
(dostęp do innych instancji PostgreSQL), \texttt{mysql\_fdw}, \texttt{mongo\_fdw},
\texttt{redis\_fdw}, \texttt{file\_fdw} (pliki CSV), \texttt{jdbc\_fdw}
oraz wrappery do S3, Kafki i systemów plików Hadoop. Do integracji
strumieniowej (CDC -- \textit{Change Data Capture}) wykorzystuje się
\textbf{Debezium} konsumujący replikację logiczną i publikujący zmiany
na Apache Kafka.

W chmurze PostgreSQL jest dostępny w modelu zarządzanym na wszystkich
trzech dużych platformach: AWS (RDS for PostgreSQL, Aurora PostgreSQL-Compatible),
Google Cloud (Cloud SQL, AlloyDB), Azure (Azure Database for PostgreSQL --
w wersjach Flexible Server oraz Hyperscale/Citus).

\subsection{Obszary biznesowe}

Dzięki ścisłej zgodności ze standardem, bogatemu katalogowi typów danych
i dojrzałości silnika, PostgreSQL jest uniwersalnym wyborem w większości
scenariuszy OLTP oraz HTAP (\textit{Hybrid Transactional/Analytical}).
Szczególnie mocną pozycję ma w następujących obszarach:

\begin{itemize}
    \item \textbf{Sektor finansowy} -- ACID, izolacja serializowalna, wsparcie
    dla typów \texttt{NUMERIC} o dowolnej precyzji (bez błędów
    zmiennoprzecinkowych), audyt przez pgaudit.
    \item \textbf{Administracja publiczna i instytucje} -- otwarta licencja
    eliminuje koszty licencyjne, a transparentność kodu jest istotna
    w kontekście suwerenności danych.
    \item \textbf{Systemy geograficzne (GIS)} -- rozszerzenie \textbf{PostGIS}
    jest de facto standardem branżowym dla operacji geoprzestrzennych,
    wykorzystywanym przez OpenStreetMap, systemy katastralne i logistyczne.
    \item \textbf{E-commerce i systemy transakcyjne} -- Shopify, Magento,
    WooCommerce natywnie wspierają PostgreSQL jako backend.
    \item \textbf{Szeroko rozumiany sektor technologiczny} -- Apple używa
    PostgreSQL w iCloud, Instagram jako głównego store'a grafu społecznego,
    Reddit jako backend-u platformy, podobnie TripAdvisor, GitLab i wiele
    innych SaaS-ów z obciążeniem rzędu miliardów rekordów.
    \item \textbf{Nauka i akademia} -- natywne wsparcie dla tablic, typów
    dokumentowych (JSONB), pełnotekstowego wyszukiwania i rozszerzeń
    statystycznych (MADlib).
\end{itemize}

\subsection{Zalety}

\begin{itemize}
    \item \textbf{Pełna zgodność ACID} i MVCC -- czytelnicy nie blokują
    pisarzy ani odwrotnie, co daje wysoką współbieżność bez zakleszczeń
    charakterystycznych dla modeli z blokadami.
    \item \textbf{Bogactwo typów danych} -- oprócz standardowego SQL:
    tablice, zakresy (\texttt{RANGE}), hstore, JSON/JSONB, UUID, typy
    geometryczne, sieciowe (CIDR, INET, MACADDR), pełnotekstowe
    (\texttt{tsvector}/\texttt{tsquery}), a także typy niestandardowe
    definiowane przez użytkownika.
    \item \textbf{Zaawansowane indeksowanie} -- B-tree, Hash, GIN (dla
    JSON, tablic i pełnotekstu), GiST (dla danych geometrycznych i zakresów),
    SP-GiST, BRIN (dla gigantycznych tabel czasoszeregowych), Bloom
    (przez rozszerzenie). Wsparcie dla indeksów częściowych, funkcyjnych,
    wielokolumnowych i pokrywających (\textit{covering} z \texttt{INCLUDE}).
    \item \textbf{Rozszerzalność} -- użytkownik może definiować własne
    typy, operatory, funkcje agregujące oraz pisać procedury w wielu
    językach (PL/pgSQL, PL/Python, PL/Perl, PL/V8/JavaScript, PL/Tcl).
    \item \textbf{Ekosystem rozszerzeń} -- PostGIS, TimescaleDB,
    pgvector (wektory embeddingów LLM), pg\_trgm, pg\_partman,
    pg\_stat\_statements, pgcrypto.
    \item \textbf{Dojrzałość i stabilność} -- aktywny rozwój od blisko 40 lat,
    silna społeczność, regularne wydania majorowe (raz w roku) wspierane
    przez 5 lat.
    \item \textbf{Licencja} -- liberalna, bez opłat licencyjnych i bez
    ograniczeń dla użycia komercyjnego.
\end{itemize}

\subsection{Wady i ograniczenia}

\begin{itemize}
    \item \textbf{Brak natywnego shardingu} -- skalowanie horyzontalne
    wymaga zewnętrznych narzędzi (Citus) lub shardingu na poziomie aplikacji.
    \item \textbf{Narzut MVCC} -- każdy \texttt{UPDATE} tworzy nową wersję
    wiersza zamiast modyfikować istniejącą, co generuje "śmieci"
    (\textit{dead tuples}) i wymaga regularnego \texttt{VACUUM}. Na obciążeniach
    z bardzo częstymi aktualizacjami może prowadzić do tzw.~\textit{table bloat}.
    \item \textbf{Wysoki koszt połączenia} -- każde połączenie to osobny
    proces systemowy (a nie lekki wątek jak w MySQL), z narzutem rzędu
    kilku megabajtów RAM. W systemach z tysiącami klientów konieczny
    jest pooler (\texttt{pgbouncer}).
    \item \textbf{Opóźnienie replikacji asynchronicznej} -- przy dużym
    obciążeniu zapisu replika może odstawać od master-a o kilka sekund,
    co wymaga świadomego projektowania aplikacji (\textit{read-after-write
    consistency}).
    \item \textbf{Złożona administracja} -- bogactwo parametrów (ponad
    300 ustawień w \texttt{postgresql.conf}) i mechanizmów (WAL,
    checkpoint, autovacuum, statystyki planera) wymaga znaczącej wiedzy
    operacyjnej w porównaniu do prostszych silników.
    \item \textbf{Mniej dojrzałe narzędzia diagnostyczne niż Oracle/SQL~Server} --
    mimo \texttt{pg\_stat\_statements}, \texttt{EXPLAIN (ANALYZE, BUFFERS)}
    i rozszerzeń monitorujących, brakuje zintegrowanego środowiska
    typu Oracle Enterprise Manager.
    \item \textbf{Case-sensitivity nieokarowanych identyfikatorów} -- standardowo
    PostgreSQL sprowadza nieokarowane identyfikatory do małych liter,
    co bywa źródłem pomyłek przy migracjach z MS~SQL Server.
\end{itemize}

\subsection{Udogodnienia}

PostgreSQL wyróżnia się szeregiem wygodnych funkcji zwiększających
produktywność na co dzień:

\begin{itemize}
    \item \textbf{Rozbudowany SQL} -- funkcje okienkowe (\texttt{OVER},
    \texttt{PARTITION BY}, \texttt{ROWS BETWEEN}), CTE (w tym rekurencyjne),
    \texttt{LATERAL JOIN}, \texttt{GROUPING SETS}/\texttt{CUBE}/\texttt{ROLLUP},
    \texttt{MERGE} (od wersji~15), \texttt{RETURNING} na \texttt{INSERT}/\texttt{UPDATE}/\texttt{DELETE}.
    \item \textbf{JSON/JSONB} -- natywne typy dokumentowe, operatory ścieżkowe
    (\texttt{->}, \texttt{->>}, \texttt{@>}, \texttt{jsonpath}) oraz indeksy
    GIN pozwalające efektywnie filtrować po dowolnych polach dokumentu.
    \item \textbf{Full-text search} -- wbudowane wyszukiwanie pełnotekstowe
    z obsługą wielu języków (w tym polskiego przez \texttt{pg\_catalog.polish}),
    tokenizacji, stop-words, steamingu i rankingu.
    \item \textbf{LISTEN/NOTIFY} -- wbudowany mechanizm pub/sub bez zewnętrznego
    brokera -- aplikacje mogą otrzymywać powiadomienia o zmianach w bazie
    w czasie rzeczywistym.
    \item \textbf{Narzędzia CLI/GUI} -- \texttt{psql} to jedno z najbogatszych
    narzędzi CLI (wsparcie dla \texttt{\textbackslash d}, \texttt{\textbackslash timing},
    \texttt{\textbackslash watch}, formatowanie wyników, podłączanie skryptów).
    GUI obejmuje pgAdmin, DBeaver, DataGrip, pg\_activity.
    \item \textbf{EXPLAIN (ANALYZE, BUFFERS, FORMAT JSON)} -- jedno
    z najlepszych narzędzi do analizy planów wykonania w branży, pozwalające
    diagnozować nie tylko typ operacji, ale też liczbę wierszy szacowanych
    vs rzeczywistych, wykorzystanie buforów i rzeczywisty czas wykonania
    per węzeł planu.
    \item \textbf{Generated columns i partial indexes} -- kolumny obliczane
    automatycznie (od wersji~12) oraz indeksy częściowe obejmujące tylko
    wybrany podzbiór wierszy (\texttt{WHERE is\_active = TRUE}) --
    zmniejszają koszt zapisu i rozmiar indeksu.
\end{itemize}
